\section{Discussion and conclusions}
\label{sec:conclusiones}

\subsection{Discussion}

The objective was to classify each station-day as an exceedance or not of
the MDA8 at 51 ppb from same-day conditions, with an AUC of at least 0.75,
and to identify the latent conditions that explain that classification. It
is met: the three classifiers exceed the threshold, and the interdependence
techniques produce three axes with a physical interpretation that is stable
outside the fitting period.

Ozone exceedance in the MMA is, above all, a seasonal and regional
phenomenon. Season and zone are the predictors with the largest effects in
the station-day models ---spring triples the odds of winter--- and
photochemical forcing, the axis of temperature, radiation and dryness, is the
dominant continuous factor. That forcing weighs more than the precursor load
is consistent with the chemistry of ozone as a secondary pollutant: the
ingredients are necessary, but it is the available energy that sets the rate
of the reactions. Ventilation is the only protective mechanism measured, and
FA shows that it is wind, not pressure, that defines this axis.

Against that seasonal background there is a spatial signal that the measured
meteorology does not explain. For the same season and conditions, the
Northwest and Southwest almost quadruple the odds of the Northeast, and a
station's exceedance rate correlates at 0.83 with its elevation. The
prevailing easterly wind carries the precursors toward the foot of the
mountains, where the orography blocks their exit; the west--east ordering is
stable across the four seasons. The zone contrasts, however, rest on 15
spatial units, and their precision is what 15 units can provide.

Temporal dependence is the third finding. Day-to-day autocorrelation survives
the seasonal decomposition in the four metropolitan-area series, and a single
lag of the response captures almost all of that dependence: having exceeded
yesterday almost quadruples the odds of exceeding today, controlling for
factors and calendar. The autoregressive model distinguishes days within a
season, which climatology cannot do.

Two limitations bound all of the above. The first is the year effect: 2024
and 2025 have an exceedance level that the factors do not explain, and no
model generalizes as well to 2021 as to 2024. With these data it cannot be
decided whether this is a change in emissions, in instrumentation, or two
meteorologically atypical years. The second is the absence of volatile
organic compounds: without them, the axes are described as consistent with
combustion, ventilation or photochemistry, and no statement about the regime
that limits ozone formation can be answered.

\subsection{Conclusions}

\begin{enumerate}[label=\roman*.]
  \item The three classifiers reach AUCs from 0.846 to 0.898 in 2025 and
    exceed the performance threshold. The logistic model and the
    discriminant are equivalent in practice, so the interpretation of the
    odds ratios does not depend on the normality assumption.
  \item Exceedance is mainly seasonal and regional: season and zone weigh
    more than any continuous factor, and among these photochemical forcing
    (2.78 per standard deviation) dominates primary combustion (1.92);
    ventilation protects (0.90).
  \item There is a spatial signal not explained by the measured meteorology:
    the west of the MMA almost quadruples the odds of the Northeast, in
    correspondence with elevation and the prevailing wind.
  \item The working calendar has little influence; only Sunday has a clear,
    small effect, and the Saturday effect disappears when standard errors
    are corrected for dependence.
  \item Day-to-day dependence is genuine and an AR(1) model exploits it: it
    outperforms climatology and persistence with an expected generalization
    of 0.878.
\end{enumerate}

\subsection{Recommendations}

For air quality management, the autoregressive model on the
metropolitan-area series can operate as an exceedance \emph{nowcasting}
system with same-day measurements, with a decision threshold that favors
sensitivity, because the cost of failing to warn of a bad day exceeds that of
one alert too many. The Northwest and Southwest zones deserve specific
monitoring and control measures, since their risk exceeds what their
meteorological conditions explain. And adding volatile organic compound
measurements to the SIMA network is the only way to move from associations
to mechanism.

For the analysis, three methodological improvements are identified:
estimating the Box-Cox parameter on the fitting data only, to close the only
remaining leak; including a moving-average term on the residuals, justified by
the second-order residual of the autoregressive model; and extending the
model to forecasting by lagging the factors, with its own validation.

\subsection{Open questions}

\begin{itemize}
  \item What produced the 2024 step? Distinguishing between a change in
    emissions, in instrumentation, and atypical years requires emission
    inventories and calibration logs that this work did not have.
  \item What characteristic of the west explains the excess risk that the
    zone retains after controlling for meteorology? Land use, distance to
    industrial sources and orographic recirculation are the candidates.
  \item How much of the season effect is photochemistry and how much is
    wind regime? The factors capture it in part, but season retains an
    effect of its own that suggests an unmeasured seasonal variable.
  \item Up to what horizon is forecasting useful? The 1.3-day mixing time of
    the binary chain suggests that beyond two or three days the prediction
    converges to the monthly climatology.
\end{itemize}
